\chapter{Unit 3: Introduction to Number System and Logic Gates}

\section{Section 1: Chapter-Wise Multiple-Choice Questions (60 MCQs)}

\subsection*{Part I: Foundational Level (Questions 1 to 20)}
\mcqitem{1}{What is the decimal equivalent of binary $1010_2$?}{$14_{10}$}{$8_{10}$}{$12_{10}$}{$10_{10}$}
\mcqitem{2}{What is the binary equivalent of decimal $13_{10}$?}{$1011_2$}{$1100_2$}{$1110_2$}{$1101_2$}
\mcqitem{3}{What is the Gray code equivalent of binary $1011_2$?}{$1010$}{$1110$}{$1001$}{$1101$}
\mcqitem{4}{What is the binary equivalent of Gray code $1101$?}{$1110$}{$1100$}{$1001$}{$1011$}
\mcqitem{5}{What is the Excess-3 code for decimal digit $4$?}{$0011$}{$1000$}{$0111$}{$0100$}
\mcqitem{6}{Which decimal digit is represented by Excess-3 code $1000$?}{$7$}{$4$}{$5$}{$6$}
\mcqitem{7}{Which BCD code represents decimal digit $9$?}{$1001$}{$1010$}{$1000$}{$1111$}
\mcqitem{8}{How is decimal number $25$ represented in BCD?}{$0010\,0101$}{$0101\,0010$}{$0010\,0110$}{$0011\,0101$}
\mcqitem{9}{What is the 1's complement of binary number $1010$?}{$1111$}{$0101$}{$0110$}{$1001$}
\mcqitem{10}{What is the 2's complement of the 4-bit binary number $0101$?}{$1110$}{$1100$}{$1011$}{$1010$}
\mcqitem{11}{What is the result of binary addition $101_2 + 011_2$?}{$110_2$}{$100_2$}{$1000_2$}{$111_2$}
\mcqitem{12}{What is the result of binary subtraction $1101_2 - 0101_2$?}{$1001_2$}{$1010_2$}{$1000_2$}{$0110_2$}
\mcqitem{13}{In 2's complement subtraction, what is added to the minuend?}{The 2's complement of subtrahend}{The 1's complement}{The sign bit}{The parity bit}
\mcqitem{14}{Using 4-bit 2's complement, what is the result of $0110_2 - 0011_2$?}{$1001_2$}{$1110_2$}{$0011_2$}{$0100_2$}
\mcqitem{15}{Which logic gate produces an output of 1 only when all inputs are 1?}{AND gate}{XOR gate}{OR gate}{NOT gate}
\mcqitem{16}{What is the output of a NOT gate when its input is 0?}{$0$}{$1$}{The same input}{Undefined}
\mcqitem{17}{According to Boolean identity law, what is $A + 0$?}{$A$}{$0$}{$1$}{$\bar{A}$}
\mcqitem{18}{What does SOP stand for in Boolean algebra?}{Sum of Products}{Sequence of Products}{Sum of Powers}{Set of Parameters}
\mcqitem{19}{Which expression is in Product of Sums (POS) form?}{$A+B+C$}{$AB+CD$}{$ABC+DEF$}{$(A+B)(C+D)$}
\mcqitem{20}{What is the main purpose of a Karnaugh map?}{To convert decimal numbers}{To generate clock pulses}{To simplify Boolean expressions}{To store binary data}

\subsection*{Part II: Intermediate Analytical Level (Questions 21 to 40)}
\mcqitem{21}{What is the binary equivalent of decimal $173_{10}$?}{$11001101_2$}{$10101101_2$}{$10101011_2$}{$10110101_2$}
\mcqitem{22}{Convert hexadecimal $3\text{A}7_{16}$ to decimal:}{$947_{10}$}{$917_{10}$}{$935_{10}$}{$927_{10}$}
\mcqitem{23}{What is the binary equivalent of Gray code $1101$?}{$1010$}{$1111$}{$1001$}{$1011$}
\mcqitem{24}{Convert binary $10110_2$ into Gray code:}{$11101$}{$10101$}{$11111$}{$10011$}
\mcqitem{25}{What is the Excess-3 representation of decimal $59$?}{$10011000$}{$01011001$}{$10001100$}{$01101001$}
\mcqitem{26}{The BCD code $1001\,0110$ represents which decimal number?}{$86$}{$96$}{$150$}{$106$}
\mcqitem{27}{What is the 1's complement of 8-bit binary $01011010$?}{$10100101$}{$11011010$}{$10101001$}{$01000101$}
\mcqitem{28}{Result of binary addition $1011_2 + 1101_2$:}{$10000_2$}{$11000_2$}{$10110_2$}{$11100_2$}
\mcqitem{29}{Using 5-bit 2's complement arithmetic, what is $10110_2 - 01001_2$?}{$00101_2$}{$01101_2$}{$10011_2$}{$11101_2$}
\mcqitem{30}{Using 5-bit signed 2's complement, what is $01101_2 + 00111_2$?}{$01010_2$}{$10100_2$}{$11100_2$}{$10010_2$}
\mcqitem{31}{A NAND gate has inputs $A=1$ and $B=1$. Output is:}{$1$}{$0$}{$A$}{$B$}
\mcqitem{32}{Which gate produces output 1 only when two inputs are different?}{XOR gate}{OR gate}{AND gate}{XNOR gate}
\mcqitem{33}{Simplify Boolean expression $A + AB$:}{$A+B$}{$AB$}{$A$}{$B$}
\mcqitem{34}{Using De Morgan's theorem, simplify $(A+B)'$:}{$AB$}{$(AB)'$}{$A'B'$}{$A'+B'$}
\mcqitem{35}{Which of the following is in standard Sum of Products (SOP) form?}{$(A+B')+C$}{$A(B+C')$}{$A'B + AC + BC'$}{$(A+B)(C+D)$}
\mcqitem{36}{Which is in Product of Sums (POS) form?}{$(A+B')(A'+C)$}{$AB(C+D)$}{$A(B+C)$}{$AB+A'C$}
\mcqitem{37}{Using 3-variable K-map, simplify $F(A,B,C) = \sum m(1,3,5,7)$:}{$A$}{$A+B$}{$C$}{$B$}
\mcqitem{38}{Using 4-variable K-map for $F = \sum m(0,2,8,10)$:}{$B'C'$}{$A'D'$}{$BD'$}{$B'D'$}
\mcqitem{39}{Using 4-variable K-map for $F = \sum m(1,3,9,11)$:}{$A'D$}{$B'C'$}{$BC'$}{$B'D$}
\mcqitem{40}{Simplify Boolean expression $(A+B)(A+C)$:}{$A+BC$}{$ABC$}{$A+B+C$}{$AB+AC$}

\subsection*{Part III: Advanced Mastery Level (Questions 41 to 60)}
\mcqitem{41}{Convert hexadecimal $(3\text{A}7.\text{C})_{16}$ to octal:}{$(1647.3)_8$}{$(1647.6)_8$}{$(1657.4)_8$}{$(1727.6)_8$}
\mcqitem{42}{Which base-7 representation is exactly equivalent to $(231.2)_4$?}{$(63.\bar{4})_7$}{$(64.3)_7$}{$(63.\bar{3})_7$}{$(63.3)_7$}
\mcqitem{43}{6-bit Gray $101101$ received with 3rd bit from left inverted. Decimal value represented by corrupted Gray word:}{$59$}{$55$}{$53$}{$57$}
\mcqitem{44}{What is binary equivalent of reflected Gray $11010111$?}{$11001010$}{$10010110$}{$10110010$}{$10011010$}
\mcqitem{45}{Two Excess-3 numbers $1001\,0100$ and $0111\,1010$. Excess-3 representation of their decimal sum:}{$0100\,0011\,1011$}{$0100\,0011\,1000$}{$0001\,0000\,1000$}{$0100\,0110\,1011$}
\mcqitem{46}{Two BCD digits $1000$ and $0111$ added with carry-in 1. Output carry and BCD sum digit after correction:}{Carry 0, digit $0110$}{Carry 1, digit $0110$}{Carry 0, digit $1001$}{Carry 1, digit $0000$}
\mcqitem{47}{Using 8-bit one's complement, what bit pattern results from $37 - 92$ before interpreting sign?}{$11001001$}{$11001000$}{$10100111$}{$00110111$}
\mcqitem{48}{Fixed-point octal register (6 integer, 2 fraction). 8's complement of $(000527.40)_8$:}{$(777250.40)_8$}{$(777470.40)_8$}{$(777251.40)_8$}{$(777250.37)_8$}
\mcqitem{49}{Evaluate unsigned binary product $(110101)_2 \times (1011)_2$:}{$(1000111111)_2$}{$(1010000111)_2$}{$(1001000111)_2$}{$(1001010111)_2$}
\mcqitem{50}{Unsigned binary $(101101101)_2$ divided by $(1101)_2$. Quotient and remainder:}{Quotient $11101$, rem 0}{Quotient $11100$, rem 11}{Quotient $11011$, rem 10}{Quotient $11100$, rem 1}
\mcqitem{51}{8-bit two's complement arithmetic for $45 - (-83)$:}{$01111111$, overflow}{$10000000$, overflow}{$00000000$, no overflow}{$10000000$, no overflow}
\mcqitem{52}{6-bit two's complement processor evaluates $(-27) - (+14)$:}{$101001$, overflow}{$010111$, no overflow}{$100111$, no overflow}{$010111$, overflow}
\mcqitem{53}{NAND logic: $X=A\,\text{NAND}\,B, Y=A\,\text{NAND}\,X, Z=B\,\text{NAND}\,X, F=Y\,\text{NAND}\,Z$. Function $F$:}{$AB$}{$A \oplus B$}{$A \odot B$}{$A+B$}
\mcqitem{54}{NOR logic: $X=A\,\text{NOR}\,B, Y=A\,\text{NOR}\,X, Z=B\,\text{NOR}\,X, F=Y\,\text{NOR}\,Z$. Function $F$:}{$A \oplus B$}{$A \odot B$}{$A+B$}{$(AB)'$}
\mcqitem{55}{Simplify $F = (A+B)(A'+C)(B+C)$ using Boolean algebra:}{$(A+B)(A'+C)$}{$(A+C)(A'+B)$}{$(A+B)(B+C)$}{$(A'+C)(B+C)$}
\mcqitem{56}{Which expression is equivalent to $F = A'B + AC + BC$?}{$AC+BC$}{$A'B+BC$}{$A'B+AC$}{$A'B+A'C$}
\mcqitem{57}{Canonical POS representation of $F(A,B,C,D) = \sum m(0,2,5,7,8,10,13,15)$:}{$\prod M(1,3,4,6,9,11,12,14)$}{$\prod M(0,2,5,7,8,10,13,15)$}{$\prod M(1,3,5,7,9,11,13,15)$}{$\prod M(2,4,6,8,10,12,14,15)$}
\mcqitem{58}{Convert $F = (A+B')(B+C)$ to canonical SOP form for $A,B,C$:}{$\sum m(0,2,3,4)$}{$\sum m(1,3,5,7)$}{$\sum m(1,5,6,7)$}{$\sum m(2,4,6,7)$}
\mcqitem{59}{4-variable K-map: $F = \sum m(1,3,9,11,12,14) + d(0,2,8,10)$ in SOP form:}{$B'+AD$}{$B'+A'D'$}{$B'D+AD'$}{$B'+AD'$}
\mcqitem{60}{Minimal 2-level SOP $F = \sum m(4,5,6,7,10,11,14,15) = A'B + AC$ has static-1 hazard. Hazard-free expression:}{$A'B + AC + BC$}{$A'B + AC + A'C$}{$A'B + AC + B'C$}{$A'B + AC + AB$}

\newpage
\section{Section 2: Complete Answer Key \& Technical Rationales}

\subsection*{Part A: Questions 1 to 30}
\begin{table}[htbp]
\centering
\caption{\textbf{Answer Key with Rationales (Questions 1 to 30)}}
\vspace{2mm}
\small
\begin{tabularx}{\textwidth}{l c X l c X}
\toprule
\textbf{Q\#} & \textbf{Ans} & \textbf{Technical Rationale} & \textbf{Q\#} & \textbf{Ans} & \textbf{Technical Rationale} \\
\midrule
1 & \textbf{D} & $1010_2 = 1\times 8 + 0\times 4 + 1\times 2 + 0 = 10_{10}$. & 16 & \textbf{B} & NOT inverts logic level: $\bar{0} = 1$. \\
2 & \textbf{D} & $13 = 8 + 4 + 1 = 1101_2$. & 17 & \textbf{A} & $A + 0 = A$. \\
3 & \textbf{B} & $G_3=1, G_2=1\oplus 0=1, G_1=0\oplus 1=1, G_0=1\oplus 1=0 \implies 1110$. & 18 & \textbf{A} & SOP = Sum of Products (OR of AND terms). \\
4 & \textbf{C} & $B_3=1, B_2=1\oplus 1=0, B_1=0\oplus 0=0, B_0=0\oplus 1=1 \implies 1001_2$. & 19 & \textbf{D} & POS is an AND of OR clauses. \\
5 & \textbf{C} & $4 + 3 = 7 = 0111_2$. & 20 & \textbf{C} & K-Maps graphically minimize Boolean switching functions. \\
6 & \textbf{C} & $1000_2 = 8; 8 - 3 = 5$. & 21 & \textbf{B} & $173 = 128 + 32 + 8 + 4 + 1 = 10101101_2$. \\
7 & \textbf{A} & 8421 code for 9 is $1001_2$. & 22 & \textbf{C} & $3\times 256 + 10\times 16 + 7 = 768 + 160 + 7 = 935_{10}$. \\
8 & \textbf{A} & $2 = 0010_2, 5 = 0101_2$. & 23 & \textbf{C} & $B_3=1, B_2=1\oplus 1=0, B_1=0\oplus 0=0, B_0=0\oplus 1=1 \implies 1001_2$. \\
9 & \textbf{B} & Invert all bits: $1010 \to 0101$. & 24 & \textbf{A} & $G_4=1, G_3=1\oplus 0=1, G_2=0\oplus 1=1, G_1=1\oplus 1=0, G_0=1\oplus 0=1 \implies 11101$. \\
10 & \textbf{C} & 1's comp $= 1010; 1010 + 1 = 1011$. & 25 & \textbf{C} & $5+3=8=1000_2, 9+3=12=1100_2 \implies 1000\,1100$. \\
11 & \textbf{C} & $5 + 3 = 8 = 1000_2$. & 26 & \textbf{B} & $1001_2 = 9, 0110_2 = 6 \implies 96$. \\
12 & \textbf{C} & $13 - 5 = 8 = 1000_2$. & 27 & \textbf{A} & Bitwise NOT: $01011010 \to 10100101$. \\
13 & \textbf{A} & $A - B = A + (\text{2's comp of } B)$. & 28 & \textbf{B} & $11 + 13 = 24 = 11000_2$. \\
14 & \textbf{C} & $6 - 3 = 3 = 0011_2$. & 29 & \textbf{B} & $22 - 9 = 13 = 01101_2$. \\
15 & \textbf{A} & Standard definition of AND gate. & 30 & \textbf{B} & $+13 + (+7) = +20$. In 5 bits, range is $[-16,+15]$; $+20 = 10100_2$ (overflow occurs, stored pattern $10100$). \\
\bottomrule
\end{tabularx}
\end{table}

\subsection*{Part B: Questions 31 to 60}
\begin{table}[htbp]
\centering
\caption{\textbf{Answer Key with Rationales (Questions 31 to 60)}}
\vspace{2mm}
\small
\begin{tabularx}{\textwidth}{l c X l c X}
\toprule
\textbf{Q\#} & \textbf{Ans} & \textbf{Technical Rationale} & \textbf{Q\#} & \textbf{Ans} & \textbf{Technical Rationale} \\
\midrule
31 & \textbf{B} & $\overline{1 \cdot 1} = 0$. & 46 & \textbf{B} & $8 + 7 + 1 = 16_{10}$. Binary sum $= 10000_2$. Add $0110_2 \implies 1\,0110_2$ (Carry 1, digit 6). \\
32 & \textbf{A} & XOR is the difference/inequality detector ($A \neq B$). & 47 & \textbf{A} & $37 = 00100101_2; 92 = 01011100_2$. 1's comp of 92 is $10100011_2$. Sum $= 00100101 + 10100011 = 11001000_2$. \\
33 & \textbf{C} & $A(1+B) = A(1) = A$ (Absorption law). & 48 & \textbf{A} & 7's complement $= (777250.37)_8$. Add 1 at LSB $(0.01) \implies (777250.40)_8$. \\
34 & \textbf{C} & $\overline{A+B} = \bar{A}\bar{B}$. & 49 & \textbf{C} & $53 \times 11 = 583 = 1001000111_2$. \\
35 & \textbf{C} & Sum (OR) of product (AND) terms. & 50 & \textbf{D} & $365 / 13 = 28$ rem 1. $28 = 11100_2$, rem $1_2$. \\
36 & \textbf{A} & Product (AND) of sum (OR) terms. & 51 & \textbf{B} & $45 - (-83) = +128$. 8-bit max is $+127$. Result wraps to $10000000_2$ with signed overflow flag set. \\
37 & \textbf{C} & Minterms $1,3,5,7$ correspond to $C=1$ independent of $A,B$. & 52 & \textbf{D} & $-27 - 14 = -41$. 6-bit range is $[-32, +31]$. Signed underflow occurs; stored bits $010111_2$ with overflow flag. \\
38 & \textbf{D} & Corner cells $0,2,8,10$ form a quad $= \bar{B}\bar{D}$. & 53 & \textbf{B} & This is the classic 4-NAND implementation of XOR ($A \oplus B$). \\
39 & \textbf{D} & Cells in columns $01, 11$ for rows $00, 10$ form a quad $= \bar{B}D$. & 54 & \textbf{B} & Dual 4-NOR circuit implements XNOR ($A \odot B$). \\
40 & \textbf{A} & $(A+B)(A+C) = AA + AC + AB + BC = A(1+C+B) + BC = A + BC$ (Distributive law). & 55 & \textbf{A} & By consensus theorem in POS form, $(B+C)$ is redundant consensus term. \\
41 & \textbf{B} & $0011\,1010\,0111.1100_2 \to (1647.6)_8$. & 56 & \textbf{C} & By consensus theorem, $BC$ is redundant consensus term. \\
42 & \textbf{A} & $(231.2)_4 = 2(16)+3(4)+1+2/4 = 45.5_{10}$. $45_{10} = (63)_7$; $0.5_{10} = (0.\bar{3}3..)_7$ or $(63.\bar{3})_7$. & 57 & \textbf{A} & Maxterms are the complement set of minterm indices: $\{1,3,4,6,9,11,12,14\}$. \\
43 & \textbf{B} & Corrupted: $100101$. Binary: $B_5=1, B_4=1, B_3=1, B_2=0, B_1=0, B_0=1 \implies 111001_2 = 57_{10}$ (or 55 depending on index). & 58 & \textbf{C} & $(A+\bar{B})(B+C) = AB + AC + \bar{B}C = \sum m(1, 5, 6, 7)$. \\
44 & \textbf{D} & XOR sequentially: $10011010_2$. & 59 & \textbf{D} & Octet covering rows $00, 10$ gives $\bar{B}$; quad covering $12, 14, 8, 10$ gives $A\bar{D}$. \\
45 & \textbf{A} & Decimals: $(9-3)(4-3) = 61$; $(7-3)(10-3) = 47$. Sum $= 61+47=108$. XS-3 of $1,0,8$ is $0100\,0011\,1011$. & 60 & \textbf{A} & Adding the consensus term $BC$ bridges adjacent groups and eliminates the static-1 hazard. \\
\bottomrule
\end{tabularx}
\end{table}

\newpage
\section{Section 3: Comprehensive Theory Questions \& Worked Solutions}
\begin{theoryq}{3.1}
State and prove De Morgan's Laws of Boolean Algebra using truth tables.
\end{theoryq}

\begin{theorysol}
1. \textbf{First Law:} $\overline{A + B} = \bar{A} \cdot \bar{B}$.\\ 
2. \textbf{Second Law:} $\overline{A \cdot B} = \bar{A} + \bar{B}$.\\ 
Proof by perfect induction (truth table verification for all 4 states shows identical columns for LHS and RHS).
\end{theorysol}

\vspace{3mm}

\begin{theoryq}{3.2}
Explain the concept of Static-1 Hazards in combinational circuits and show how to design hazard-free logic using Karnaugh Maps.
\end{theoryq}

\begin{theorysol}
A static-1 hazard is a transient 0 glitch that occurs during input transitions between adjacent 1s located in different prime implicant groups due to unequal gate propagation delays. It is eliminated by including the redundant consensus prime implicant covering the boundary cells.
\end{theorysol}

\vspace{3mm}

